
\documentclass{article}
\usepackage{colortbl}
\usepackage{makecell}
\usepackage{multirow}
\usepackage{supertabular}

\begin{document}

\newcounter{utterance}

\centering \large Interaction Transcript for game `referencegame', experiment `random\_grids', episode 0 with NEW\_Reflect\_Agent.
\vspace{24pt}

{ \footnotesize  \setcounter{utterance}{1}
\setlength{\tabcolsep}{0pt}
\begin{supertabular}{c@{$\;$}|p{.15\linewidth}@{}p{.15\linewidth}p{.15\linewidth}p{.15\linewidth}p{.15\linewidth}p{.15\linewidth}}
    \# & $\;$A & \multicolumn{4}{c}{Game Master} & $\;\:$B\\
    \hline

    \theutterance \stepcounter{utterance}  
    & & \multicolumn{4}{p{0.6\linewidth}}{
        \cellcolor[rgb]{0.9,0.9,0.9}{
            \makecell[{{p{\linewidth}}}]{
                \texttt{\tiny{[P1$\langle$GM]}}
                \texttt{You are given three grids, where each of them is 5 by 5 in size.} \\
\texttt{Grids have empty cells marked with "▢" and filled cells marked with "X".} \\
\texttt{Your task is to generate a referring expression that best describes the target grid while distinguishing it from the two other distractor grids.} \\
\texttt{The first grid is the target grid, and the following two grids are the distractors.} \\
\\ 
\texttt{Target grid:} \\
\\ 
\texttt{X X □ X □} \\
\texttt{□ X X □ X} \\
\texttt{X □ □ X □} \\
\texttt{□ X □ □ X} \\
\texttt{X □ X □ X} \\
\\ 
\texttt{Distractor grid 1:} \\
\\ 
\texttt{X X □ □ X} \\
\texttt{□ X X □ X} \\
\texttt{X □ X X □} \\
\texttt{□ X □ X X} \\
\texttt{X □ □ X X} \\
\\ 
\texttt{Distractor grid 2:} \\
\\ 
\texttt{X □ X X □} \\
\texttt{□ X □ X X} \\
\texttt{X X □ □ X} \\
\texttt{□ X X □ X} \\
\texttt{X □ X □ □} \\
\\ 
\texttt{Instruction: Describe the target grid.} \\
\texttt{Generate the referring expression starting with the tag "Expression: " for the given target grid. Omit any other text.} \\
            }
        }
    }
    & \\ \\

    \theutterance \stepcounter{utterance}  
    & \multicolumn{4}{p{0.45\linewidth}}{
        \cellcolor[rgb]{0.95,0.95,0.95}{
            \makecell[{{p{\linewidth}}}]{
                \texttt{\tiny{[P1$|$internal]}}
                \texttt{Reflections from previous episodes:} \\
\texttt{{-}} \\
\texttt{ Expression: A vertical stripe of X's in the middle three columns, with empty cells on the leftmost and rightmost columns.} \\
\\ 
\texttt{Reflection: Succeeded. I identified the key distinguishing feature: the target has a solid vertical block of X's occupying columns 2–4, while both distractors have the block shifted left or right, leaving two empty columns on one side. I described it clearly and concisely. Next time, I would use the same approach—focus on the shape and position of the filled region relative to the grid edges.} \\
\texttt{{-}} \\
\texttt{ I succeeded. The target grid has three filled columns on the left and two empty columns on the right, while the distractors have filled columns shifted or with gaps. I correctly described it as "three leftmost columns filled, two rightmost columns empty." This was precise and distinguished it clearly. Next time, I would use the same approach: identify the unique pattern of filled vs. empty columns/rows compared to distractors.} \\
\texttt{{-}} \\
\texttt{ I succeeded. The target grid had three filled columns on the left and two empty columns on the right, forming a solid vertical block. The distractors had filled columns shifted or with gaps. My expression "Three leftmost columns completely filled, two rightmost columns empty" clearly distinguished it. Next time, I would use the same approach: identify the unique contiguous pattern and describe its position and shape precisely.} \\
            }
        }
    }
    & & \\ \\

    \theutterance \stepcounter{utterance}  
    & \multicolumn{4}{p{0.45\linewidth}}{
        \cellcolor[rgb]{0.95,0.95,0.95}{
            \makecell[{{p{\linewidth}}}]{
                \texttt{\tiny{[P1$|$internal]}}
                \texttt{{[}attempt 1 – valid{]}} \\
\texttt{Expression: A vertical stripe of X's in the middle three columns, with empty cells on the leftmost and rightmost columns.} \\
            }
        }
    }
    & & \\ \\

    \theutterance \stepcounter{utterance}  
    & \multicolumn{4}{p{0.6\linewidth}}{
        \cellcolor[rgb]{0.9,0.9,0.9}{
            \makecell[{{p{\linewidth}}}]{
                \texttt{\tiny{[P1$\rangle$GM]}}
                \texttt{Expression: A vertical stripe of X's in the middle three columns, with empty cells on the leftmost and rightmost columns.} \\
            }
        }
    }
    & & \\ \\

    \theutterance \stepcounter{utterance}  
    & & & \multicolumn{2}{p{0.3\linewidth}}{
        \cellcolor[rgb]{0.9,0.9,0.9}{
            \makecell[{{p{\linewidth}}}]{
                \texttt{\tiny{[GM$|$GM]}}
                \texttt{Expression: A vertical stripe of X's in the middle three columns, with empty cells on the leftmost and rightmost columns.} \\
            }
        }
    }
    & & \\ \\

    \theutterance \stepcounter{utterance}  
    & & \multicolumn{4}{p{0.6\linewidth}}{
        \cellcolor[rgb]{0.9,0.9,0.9}{
            \makecell[{{p{\linewidth}}}]{
                \texttt{\tiny{[GM$\rangle$P2]}}
                \texttt{You are given three grids, where each of them is 5 by 5 in size.} \\
\texttt{Grids have empty cells marked with "▢" and filled cells marked with "X".} \\
\texttt{You are also given a referring expression that describes one of the given grids.} \\
\texttt{Your task is to select the grid that matches the given referring expression.} \\
\texttt{Generate only the number (in text) of the grid that matches the given expression by selecting first, second, or third.} \\
\\ 
\texttt{First:} \\
\\ 
\texttt{X □ X X □} \\
\texttt{□ X □ X X} \\
\texttt{X X □ □ X} \\
\texttt{□ X X □ X} \\
\texttt{X □ X □ □} \\
\\ 
\texttt{Second:} \\
\\ 
\texttt{X X □ □ X} \\
\texttt{□ X X □ X} \\
\texttt{X □ X X □} \\
\texttt{□ X □ X X} \\
\texttt{X □ □ X X} \\
\\ 
\texttt{Third:} \\
\\ 
\texttt{X X □ X □} \\
\texttt{□ X X □ X} \\
\texttt{X □ □ X □} \\
\texttt{□ X □ □ X} \\
\texttt{X □ X □ X} \\
\\ 
\texttt{Expression: A vertical stripe of X's in the middle three columns, with empty cells on the leftmost and rightmost columns.} \\
\texttt{Question: Which grid does the expression refer to?} \\
\texttt{Start with the tag "Answer: ", followed by your selection. Omit any other text.} \\
            }
        }
    }
    & \\ \\

    \theutterance \stepcounter{utterance}  
    & & & \multicolumn{4}{p{0.45\linewidth}}{
        \cellcolor[rgb]{0.95,0.95,0.95}{
            \makecell[{{p{\linewidth}}}]{
                \texttt{\tiny{[P2$|$internal]}}
                \texttt{Reflections from previous episodes:} \\
\texttt{{-}} \\
\texttt{ Answer: Second} \\
\\ 
\texttt{Reflection: Succeeded. The expression described a solid vertical block of X's in the right three columns, with left two columns empty. Grid 2 matches exactly: left two columns empty, right three columns filled. I correctly parsed the description and matched it to the visual layout. Next time, I will continue to carefully align column descriptions with the grid orientation.} \\
\texttt{{-}} \\
\texttt{ Answer: Third} \\
\\ 
\texttt{Reflection: Succeeded. Correctly matched the description of "vertical block of X's in the left three columns, with empty cells in the right two columns" to the third grid. The first grid has gaps in the fourth column, and the second has X's in the fourth column, so only the third fits. Next time, I would continue to carefully compare each grid column by column against the expression.} \\
\texttt{{-}} \\
\texttt{ Failed. I misinterpreted "three leftmost columns filled" as three columns of X's starting from the left edge, but the expression means the *first three columns* are filled and the *last two* are empty. Grid 2 matches this exactly. Next time, I will map the description directly to column indices (columns 1{-}3 filled, 4{-}5 empty) without reinterpreting "leftmost" loosely.} \\
            }
        }
    }
     \\ \\

    \theutterance \stepcounter{utterance}  
    & & & \multicolumn{4}{p{0.45\linewidth}}{
        \cellcolor[rgb]{0.95,0.95,0.95}{
            \makecell[{{p{\linewidth}}}]{
                \texttt{\tiny{[P2$|$internal]}}
                \texttt{{[}attempt 1 – valid{]}} \\
\texttt{Answer: Third} \\
            }
        }
    }
     \\ \\

    \theutterance \stepcounter{utterance}  
    & & & \multicolumn{4}{p{0.6\linewidth}}{
        \cellcolor[rgb]{0.9,0.9,0.9}{
            \makecell[{{p{\linewidth}}}]{
                \texttt{\tiny{[GM$\langle$P2]}}
                \texttt{Answer: Third} \\
            }
        }
    }
     \\ \\

    \theutterance \stepcounter{utterance}  
    & & & \multicolumn{2}{p{0.3\linewidth}}{
        \cellcolor[rgb]{0.9,0.9,0.9}{
            \makecell[{{p{\linewidth}}}]{
                \texttt{\tiny{[GM$|$GM]}}
                \texttt{third} \\
            }
        }
    }
    & & \\ \\

\end{supertabular}
}

\end{document}
